% 学工智伴 v4.9.0 完整功能产品手册
% 本文件只说明功能、使用方法、数据边界和安全规则，不包含价格、购买、激活码或商业授权内容。
\documentclass[UTF8,zihao=-4,oneside]{ctexrep}

\usepackage[a4paper,left=25mm,right=25mm,top=24mm,bottom=23mm,headheight=15mm,headsep=7mm,footskip=12mm]{geometry}
\usepackage{graphicx}
\usepackage{booktabs}
\usepackage{longtable}
\usepackage{tabularx}
\usepackage{array}
\usepackage{enumitem}
\usepackage{fancyhdr}
\usepackage{titlesec}
\usepackage{caption}
\usepackage{float}
\usepackage{hyperref}
\usepackage{xcolor}
\usepackage{tcolorbox}
\usepackage{setspace}
\usepackage{microtype}
\usepackage{makecell}

% X 列使用左对齐，避免窄表格中的短句产生无意义的 Underfull hbox 警告。
\newcolumntype{Y}{>{\raggedright\arraybackslash}X}
\newcolumntype{P}[1]{>{\raggedright\arraybackslash}p{#1}}
\newcolumntype{B}[1]{>{\bfseries\raggedright\arraybackslash}p{#1}}

% 本地构建时优先使用临时目录中的 Noto Sans SC 字体；字体文件不进入仓库。
% 没有该文件时，依次使用常见的系统 CJK 字体，并在均不可用时明确报错。
\IfFileExists{../tmp/tools/NotoSansSC-VF.ttf}{%
  \setCJKmainfont[Path=../tmp/tools/,Extension=.ttf,AutoFakeBold=1.5]{NotoSansSC-VF}%
  \setCJKsansfont[Path=../tmp/tools/,Extension=.ttf,AutoFakeBold=1.5]{NotoSansSC-VF}%
  \setCJKmonofont[Path=../tmp/tools/,Extension=.ttf,AutoFakeBold=1.5]{NotoSansSC-VF}%
}{%
  \IfFontExistsTF{Noto Sans CJK SC}{%
    \setCJKmainfont{Noto Sans CJK SC}%
    \setCJKsansfont{Noto Sans CJK SC}%
    \setCJKmonofont{Noto Sans CJK SC}%
  }{%
    \IfFontExistsTF{Microsoft YaHei}{%
      \setCJKmainfont{Microsoft YaHei}%
      \setCJKsansfont{Microsoft YaHei}%
      \setCJKmonofont{Microsoft YaHei}%
    }{%
      \PackageError{product-guide}{No CJK font found}{Install Noto Sans CJK SC or Microsoft YaHei, then rebuild this guide.}%
    }%
  }%
}

\definecolor{CWBblue}{HTML}{173F7A}
\definecolor{CWBink}{HTML}{1F2937}
\definecolor{CWBgray}{HTML}{5B6573}
\definecolor{CWBline}{HTML}{D9E0EA}
\definecolor{CWBfill}{HTML}{F4F6F8}
\definecolor{CWBnote}{HTML}{EEF3FA}
\definecolor{CWBwarn}{HTML}{F7F3E8}

\hypersetup{
  colorlinks=true,
  linkcolor=CWBblue,
  urlcolor=CWBblue,
  citecolor=CWBblue,
  pdftitle={学工智伴 v4.9.0 完整功能产品手册},
  pdfauthor={学工智伴维护组},
  bookmarksopen=true
}

% hyperref/atveryend records the physical page count in the auxiliary file.
% This remains correct when the front matter uses Roman numerals.
\makeatletter
\newcommand{\cwbtotalpages}{\@ifundefined{@abspage@last}{0}{\@abspage@last}}
\makeatother

\graphicspath{{../assets/screenshots/v4.9.0/}}
% 用户手册采用统一的块状段落，不使用首行缩进，避免中英文混排时
% 出现有的段落缩进、有的段落顶格的视觉差异。
\setlength{\parindent}{0pt}
\setlength{\parskip}{0.58em}
\linespread{1.35}
\setlist[itemize]{leftmargin=2.2em,itemsep=0.25em,topsep=0.35em}
\setlist[enumerate]{leftmargin=2.4em,itemsep=0.25em,topsep=0.35em}
\setlist[description]{leftmargin=3.2em,labelindent=0em,itemsep=0.25em,topsep=0.35em}
\renewcommand{\arraystretch}{1.35}
\setlength{\tabcolsep}{5pt}

\titleformat{\chapter}[display]
  {\normalfont\bfseries\color{CWBink}}
  {\filleft\Large 第\thechapter 章}
  {0.6ex}
  {\LARGE\filright}
  [\vspace{0.6ex}\titlerule]
\titleformat{\section}
  {\normalfont\Large\bfseries\color{CWBblue}}
  {\thesection}{0.7em}{}
\titleformat{\subsection}
  {\normalfont\large\bfseries\color{CWBink}}
  {\thesubsection}{0.7em}{}
\titleformat{\subsubsection}
  {\normalfont\normalsize\bfseries\color{CWBgray}}
  {\thesubsubsection}{0.7em}{}
\titlespacing*{\chapter}{0pt}{0pt}{22pt}
\titlespacing*{\section}{0pt}{18pt}{8pt}
\titlespacing*{\subsection}{0pt}{12pt}{5pt}
\titlespacing*{\subsubsection}{0pt}{8pt}{3pt}

\captionsetup{font=small,labelfont=bf,labelsep=space,justification=centering}
\pagestyle{fancy}
\fancyhf{}
\fancyhead[L]{\small\color{CWBgray}学工智伴}
\fancyhead[R]{\small\color{CWBgray}完整功能产品手册}
\fancyfoot[C]{\small\color{CWBgray}第 \thepage 页 / 共 \cwbtotalpages 页}
\renewcommand{\headrulewidth}{0.4pt}
\renewcommand{\footrulewidth}{0pt}
\fancypagestyle{plain}{%
  \fancyhf{}
  \fancyfoot[C]{\small\color{CWBgray}第 \thepage 页 / 共 \cwbtotalpages 页}
  \renewcommand{\headrulewidth}{0pt}
}

\newcommand{\productname}{学工智伴}
\newcommand{\version}{v4.9.0 体验版}
\newcommand{\code}[1]{\texttt{\small #1}}
\newcommand{\term}[1]{\textbf{#1}}
\newcommand{\placeholder}[1]{\texttt{\textcolor{CWBblue}{#1}}}
\newcommand{\uiShot}[3]{%
  \begin{figure}[H]
    \centering
    \includegraphics[width=\textwidth]{#1}
    \caption{#2}
    \label{#3}
  \end{figure}%
}

\newtcolorbox{infobox}{
  colback=CWBnote,
  colframe=CWBblue,
  boxrule=0.45pt,
  arc=1.5mm,
  left=5mm,right=5mm,top=3mm,bottom=3mm,
  before skip=8pt,after skip=10pt
}
\newtcolorbox{warningbox}{
  colback=CWBwarn,
  colframe=CWBgray,
  boxrule=0.45pt,
  arc=1.5mm,
  left=5mm,right=5mm,top=3mm,bottom=3mm,
  before skip=8pt,after skip=10pt
}
\newtcolorbox{workflowbox}{
  colback=CWBfill,
  colframe=CWBline,
  boxrule=0.45pt,
  arc=1.5mm,
  left=5mm,right=5mm,top=3mm,bottom=3mm,
  before skip=8pt,after skip=10pt
}

\begin{document}

\begin{titlepage}
  \thispagestyle{empty}
  \vspace*{20mm}
  \begin{center}
    {\color{CWBblue}\rule{0.72\textwidth}{1.3pt}}\\[12mm]
    {\Huge\bfseries\color{CWBink} \productname\par}
    \vspace{7mm}
    {\LARGE\bfseries 完整功能产品手册\par}
    \vspace{8mm}
    {\large\color{CWBgray}面向辅导员日常工作的功能说明与操作指南\par}
    \vfill
    \begin{tcolorbox}[width=0.78\textwidth,colback=CWBfill,colframe=CWBline,boxrule=0.5pt,arc=1.5mm]
      \centering
      \large
      \textbf{适用版本：}\version\\[4pt]
      \textbf{文档范围：}学生台账、日常工作、业务协同、智能辅助与数据管理\\[4pt]
      \textbf{阅读方式：}先看“首次使用”，再按工作场景查阅对应模块
    \end{tcolorbox}
    \vfill
    {\color{CWBgray}\rule{0.72\textwidth}{0.4pt}}\\[5mm]
    {\large 2026 年 8 月}\par
    \vspace{4mm}
    {\small\color{CWBgray}截图使用虚构演示数据，仅用于说明界面和工作流程。}
  \end{center}
\end{titlepage}

\pagenumbering{Roman}
\chapter*{阅读说明}
\addcontentsline{toc}{chapter}{阅读说明}

这是一份给实际使用者准备的操作手册。当前对外产品名称为“学工智伴”，历史截图或旧构建物中的“一站式辅导员工作台”属于兼容名称。内容按照辅导员的日常工作顺序编排，先说明模块适合解决什么问题，再说明常用按钮、查看结果和后续跟进方式。第一次使用时，建议先用示例数据走一遍流程，再录入自己的工作内容。

\begin{infobox}
本手册把重点放在日常工作和实际操作上。涉及 AI、敏感资料、备份和跨端使用的提示会保留在对应章节，方便第一次使用时按提示完成设置。
\end{infobox}

截图中的姓名、学号、联系方式和业务记录均为虚构示例。当前候选截图保留了部分历史标题栏，用于说明界面结构，不代表对外品牌名称。涉及心理、资助、纪律、住宿、成绩和组织发展的正式结论，仍应按学校制度完成审核；工作台负责把记录、提醒和材料整理得更清楚。

\tableofcontents
\clearpage
\listoffigures
\clearpage
\pagenumbering{arabic}

\chapter{产品定位与整体工作方法}

\section{它解决的核心问题}

\productname{}是一套本地优先的高校学生工作空间。它把学生台账、谈心谈话、家校联系、任务、工作留痕、成绩与帮扶、班团组织、党团发展、住宿、查课查寝、奖惩、活动、资料、就业、科研和导入导出放在同一条可回看的工作脉络中。

它的核心价值不是把更多表格搬到网页上，而是让已经发生的一次工作能够继续被找到、被跟进、被复核：

\begin{itemize}
  \item 名单导入后，学生稳定身份、当前学号、历史学号和业务记录继续保持关联。
  \item 谈话、家长联系、任务、查课、查寝、活动和通知转化都可以生成待确认工作记录草稿。
  \item 成绩、考勤、获奖、住宿、重点关注和班委履职可以按学生、班级、学期和日期回看。
  \item 照片、证书、活动材料和业务附件进入附件仓，业务记录只保存附件 ID。
  \item 导入、批量编辑、恢复、同步和删除都尽量提供差异预览、恢复点、来源回链和失败重试。
\end{itemize}

\section{一条完整的工作链}

日常使用建议形成下面的闭环：

\begin{workflowbox}
发现对象或事项 → 进入学生/班级上下文 → 记录事实 → 安排下一步 → 生成待确认草稿 → 人工核对 → 写入正式留痕 → 在首页和时间线回看 → 备份或交接
\end{workflowbox}

例如，辅导员发现一名学生连续缺勤，可以先在考勤或查课模块记录事实，再在学生档案查看成绩和近期谈话，随后建立学业帮扶或回访任务。AI 可以帮助整理文字或解释趋势，但不能代替老师判断风险、确认处分、改变成绩或关闭预警。

\section{适用范围}

\begin{table}[H]
\centering
\caption{本说明书的使用范围}
\begin{tabularx}{\textwidth}{B{30mm} Y}
\toprule
范围 & 说明 \\
\midrule
适用版本 & v4.9.0 前瞻版，适合前瞻体验用户先用示例数据熟悉完整工作流程。 \\
适合人群 & 高校一线辅导员、年级负责人和需要整理学生工作材料的老师。 \\
使用方式 & 浏览器、单文件离线版和 Windows 桌面版均可体验；长期保存真实资料时建议优先使用桌面版。 \\
资料范围 & 学生台账、谈话联系、任务留痕、成绩帮扶、住宿考勤、组织发展、资料附件和工作分析。 \\
\bottomrule
\end{tabularx}
\end{table}

\chapter{首次启动与界面阅读}

\section{先用示例数据熟悉产品}

首次打开建议先使用虚构样例数据浏览。样例学生、任务、谈话、成绩、住宿、课表、班团、资料、科研和通知带有演示标记，只用于理解按钮和练习流程，不应当直接当作真实事实导出或上报。

推荐的第一次体验路径如下：

\begin{enumerate}
  \item 打开“今日概览”，先理解摘要条、最紧急事项和右侧上下文区。
  \item 进入“学生台账”，切换表格、卡片和照片模式，查看一名学生的档案。
  \item 在档案中打开时间线，观察谈话、成绩、住宿和工作记录如何按学生稳定身份聚合。
  \item 新建一条谈话或任务，体验“保存成功后才提示”的流程；不要把样例记录直接用于真实统计。
  \item 打开“备份与迁移”，生成一次本地恢复点，了解桌面端、浏览器端和离线 HTML 的差异。
  \item 在“AI 智能工作台”查看模型状态、用途授权和建议中心；没有配置模型时，页面会明确显示不可调用原因。
\end{enumerate}

\section{桌面三栏结构}

桌面端采用稳定的三段式工作区：左侧是模块导航，中间是当前任务，右侧是上下文辅助。右侧不是装饰面板，它用于减少反复跳页：当前学生、当前事项、今日待办、待确认草稿和当前页面可执行的快捷动作都可以在这里看到。

\begin{table}[H]
\centering
\caption{不同视口下的界面退化规则}
\begin{tabularx}{\textwidth}{B{42mm} Y}
\toprule
视口范围 & 界面行为 \\
\midrule
不小于 1280px & 左侧导航、中央工作区和右侧上下文区同时显示。 \\
1024px 至 1279px & 右侧上下文区默认收起，需要时通过按钮打开。 \\
不大于 900px & 使用导航抽屉、单列主区和底部快捷操作，不强行压缩为三栏。 \\
手机窄屏 & 表格在自己的滚动区域内滚动；卡片、筛选、保存、编辑和返回按钮保持可操作。 \\
\bottomrule
\end{tabularx}
\end{table}

\begin{figure}[H]
  \centering
  \includegraphics[width=\textwidth]{01-overview.png}
  \caption{今日工作台的三栏结构：左侧导航、中央任务区与右侧上下文区}
  \label{fig:overview}
\end{figure}

\section{页面的固定阅读顺序}

多数模块遵循同一结构：页面标题与主操作、统计摘要、筛选区、结果列表或时间线、详情和导出。主按钮只承担一个主要动作；导入、模板、导出、删除、撤销和取消使用次级按钮层级。状态标签只表达事实状态，例如“待确认”“逾期”“未记录”“已归档”，不用颜色替代文字。

手机端首次打开可以从左上角菜单进入导航抽屉。抽屉外点击、返回键或 Escape 可以关闭抽屉，焦点会回到菜单按钮。导航分组可折叠，也可以通过“搜索模块”找到被折叠的入口。

\chapter{今日工作台}

\section{首页先回答什么问题}

首页围绕“今天先做什么”组织信息。顶部摘要只显示少量关键数字，默认包括未完成任务、今日到期、逾期事项、待回访、待确认工作记录和近期开题或课题节点。数字本身可以点击进入对应模块，不需要重新设置筛选条件。

主区默认列出最紧急的五条事项，按逾期、今日、近期和待确认分组。事项过多时使用折叠和“查看全部”，不把全部待办一次性铺成长页面。折叠状态保存在当前工作区，移动端使用两行摘要或横向指标条。

\section{需要留意与系统状态}

右侧区域用于展示待回访谈话、未解除风险、待确认草稿和常用快捷动作。底部或数据概览区显示最近备份、局域网连接、离线队列、同步冲突、待确认照片和最近导入等系统状态。

首页只读取摘要和必要元数据，照片、证书原文和大图表按需加载，避免首屏读取所有二进制附件。

\begin{description}
  \item[未完成任务] 只统计当前仍需处理的任务，不把草稿误算为完成事实。
  \item[待回访] 来源是谈话记录中的回访日期和完成状态，删除谈话后不会静默删除已有回访任务。
  \item[待确认草稿] 来源记录、AI 建议或通知转化产生的候选留痕，确认前不会进入正式统计。
  \item[同步冲突] 只显示待人工选择的冲突数量，不以“最后修改时间”自动覆盖同字段内容。
  \item[备份状态] 显示最近一次真实备份时间、备份后的变更数和下一次检查时间。
\end{description}

\chapter{学生台账与学生档案}

\section{学生台账的三种阅读模式}

学生台账是所有学生业务的入口。表格适合大批量核对，卡片适合日常浏览，照片名册适合快速识别和走访。三种模式使用同一份学生数据，不会因为切换显示方式而生成副本。

\uiShot{05-students-cards.png}{学生台账卡片模式：快速浏览班级、宿舍、导师、班主任和家长关系}{fig:student-cards}

\uiShot{06-students-photos.png}{照片管理：批量导入、花名册查看、待补采统计和人工归档入口}{fig:student-photos}

\begin{figure}[H]
  \centering
  \includegraphics[width=\textwidth]{02-students-pagination-bulk.png}
  \caption{学生台账：筛选、分页、列视图和批量管理}
  \label{fig:students}
\end{figure}

\section{查询、筛选与分页}

可以按姓名、当前学号、历史学号、班级、年级、性别、学生类型、住宿状态、导师、班主任、关注类型、危机等级、预警状态和自定义字段筛选。筛选与排序可以保存为当前工作区的常用视图，列顺序、列宽、每页条数和显示模式也会保留。

批量选择时必须先确认范围：当前页，还是整个筛选结果。页面会显示已选人数、影响字段和可撤销入口。大名单使用分页和原子替换策略，不能通过直接操作底层 IndexedDB 记录来“删一行”。

\section{学生档案字段}

档案编辑页可维护以下信息：

\begin{itemize}
  \item 基础身份：稳定学生 ID、当前学号、历史学号、姓名、性别、出生信息、年级、专业、班级、培养层次和学生状态。
  \item 工作关系：导师、班主任、班级组织任职、学生类型和就业状态。
  \item 联系信息：本人电话、邮箱、QQ、家长关系、家长电话、紧急联系人和家庭地址。
  \item 居住信息：校内或校外居住类型、宿舍快照、房东电话和校外地址。
  \item 专项信息：政治面貌、退伍士兵身份、奖助勤分类、伙食补助、困难认定材料和自定义字段。
  \item 附件信息：照片、证书、活动材料和其他业务附件的附件 ID、文件名、哈希和引用关系。
\end{itemize}

家长电话默认脱敏，查看完整号码需要经过访问锁确认并留下访问审计。照片、证书和附件不会直接塞进学生记录，而是进入附件仓后由学生记录保存引用。

\begin{figure}[H]
  \centering
  \includegraphics[width=0.96\textwidth]{04-student-timeline.png}
  \caption{学生档案：基础信息、重点状态与跨模块时间线}
  \label{fig:student-timeline}
\end{figure}

\section{稳定身份与长期维护}

所有学生业务关系优先使用稳定的 \code{student\_id}。当前学号是快照，历史学号是兼容索引。学号变更时保留原学生主体，并把旧学号追加到历史记录中；谈话、成绩、住宿、奖惩、资助、就业和班团记录不会因学号修正而断开。

如果一条旧记录没有稳定 ID，可以使用当前学号或唯一历史学号提供候选，但姓名永远不自动匹配。同名、多个候选或歧义历史学号会进入“身份核对”队列，老师选择学生并确认后，系统只补写稳定关系和缺失快照，不改写原始事实。

\section{时间线与关联影响}

时间线将谈话、家校联系、成绩、预警、帮扶、住宿变动、奖惩、活动、工作留痕、党团发展和班委任职按日期聚合。它用于理解上下文，不会把学生档案修改自动传播为对历史成绩、心理、纪律、资助和预警事实的改写。

编辑学生前，页面可以显示本次变更影响的关联模块和记录数量。班级、宿舍和进行中任务等有明确快照关系的内容按业务规则同步；过程事实保持原始记录和来源，不做静默重写。

\chapter{学生导入、增量更新与批量操作}

\section{导入前的正确准备}

导入前建议从“模板中心”下载模板，保留第一行表头和字段说明，在副本上整理数据。第一次导入真实数据时先使用一个班级或一小批学生，核对数量、学号、姓名、班级、照片引用和敏感字段，再扩大范围。

系统支持 CSV 和 XLSX。表头会按规范字段名、中文显示名和学校自定义字段映射；无法识别的列不会静默丢弃，而是保留为待确认的自定义字段。

\section{合并更新规则}

默认模式是“合并更新”：来源表中没有出现的学生不会被删除；来源文件中没有出现的字段不会覆盖旧值；空白字段默认保留原值。只有明确勾选“清空空值”，某列中的空值才有清空含义。

身份匹配顺序为：

\begin{enumerate}
  \item 文件中存在且匹配的稳定 \code{student\_id}；
  \item 当前学号；
  \item 唯一历史学号；
  \item 无法唯一判断时进入人工核对；
  \item 姓名不参与自动合并。
\end{enumerate}

例如，期末只导入“学号、姓名、分数”三列时，系统会先通过学号定位已有学生，只更新成绩相关字段，培养层次、家长电话、导师和住宿等未出现在文件中的内容保持不变。课程、学期、绩点和挂科状态需要使用成绩模板导入，不应把课程化成绩当成简单学生名单字段。

\begin{figure}[H]
  \centering
  \includegraphics[width=\textwidth]{03-import-preview.png}
  \caption{导入预览：表头识别、合并模式、差异统计和冲突提示}
  \label{fig:import-preview}
\end{figure}

\section{预览、覆盖与撤销}

确认导入前会显示新增、更新、待核对、冲突和错误数量，并可展开查看逐字段差异。覆盖导入只有在明确确认源文件是完整名单后才允许执行，确认前会显示可能被归档或删除的数量，并自动生成恢复点。

写入失败时，页面保留原始文件、导入预览和重试入口，学生数据与班级历史一同回滚。看到成功提示前不要关闭窗口或刷新页面；成功后建议刷新台账并抽查几名学生。

\section{自定义字段}

“学生字段中心”用于维护学校自定义字段目录。每个字段可以定义显示名称、类型、可选值、敏感级别、是否允许导入导出和更新时间。模板中同时显示示例值、日期格式、敏感标记和清空规则。

\uiShot{50-student-fields-and-history.png}{学生字段中心：学校可以自己维护字段、选项和敏感边界}{fig:student-fields}

删除字段目录只删除字段定义，不删除学生已有值。若学校决定停用某字段，建议先导出保留记录并在制度允许的范围内归档，避免把“字段不再显示”误解成“数据已经删除”。

\section{批量编辑、归档与删除}

批量编辑支持班级、导师、班主任、家长字段、住宿类型和自定义字段。批量归档和删除必须再次确认范围，并创建恢复点。删除涉及附件或业务关联时，系统先检查引用关系；共享附件不会因为删除一条记录而被误删。

批量操作完成后可以撤销最近一次操作。撤销是针对同一工作区的恢复动作，不应替代日常备份。批量失败时恢复原列表、队列、附件引用和正在编辑的表单。

\chapter{谈心谈话、家校联系与工作留痕}

\section{谈心谈话记录}

谈话记录建议至少填写学生、日期、方式、主题、谈话要点、结果、下一步、回访日期和附件。重点关注、风险和危机字段用于工作跟进，不等同于医学诊断或专业心理结论。

谈话列表支持按学生姓名、学号、稳定 ID、主题、日期区间和跟进状态筛选。逾期回访会回到首页摘要；完成回访后标记完成，并可记录新的跟进内容。

\begin{figure}[H]
  \centering
  \includegraphics[width=\textwidth]{06-talk-crisis-schedule.png}
  \caption{谈心谈话：筛选、跟进日期、状态和学生上下文}
  \label{fig:talks}
\end{figure}

\section{会前简报与谈后草稿}

会前简报可以汇总当前学生的档案摘要、近期谈话、任务、成绩风险、住宿变动和待核对事项。它只是工作提示，不能自动生成危机结论。

谈话保存后可以生成一条待确认工作记录草稿。草稿保留来源集合、来源记录、来源哈希、来源更新时间和人工确认状态。老师可以编辑、合并、驳回或确认；来源被修改后必须重新核对，来源被删除后不能直接归档。

\uiShot{08-talks.png}{谈心谈话工作区：记录、回访日期和待跟进学生集中查看}{fig:talk-workspace}

\section{家校联系}

家长联系单独保存在家校联系记录中，与“工作通讯录”严格分离。记录包含家长关系快照、联系日期、联系方法、联系目的、沟通摘要、沟通结果、下一步和附件。

家长基础字段是学生档案的一部分，家长联系记录是一次实际发生的沟通事实。更新家长电话不会改写过去记录中的关系快照；过去联系记录保留当时的来源和日期，便于交接与复核。

\uiShot{35-family-research.png}{家校联系：家长关系、沟通摘要、结果和下一步单独留存}{fig:family-contact}

\section{统一工作留痕}

以下动作可以生成同一套待确认草稿链：

\begin{itemize}
  \item 谈心谈话和家长联系；
  \item 任务完成、回访完成和工作节点完成；
  \item 查课、随机点名和查寝；
  \item 活动或班会完成；
  \item 通知 AI 结果确认后转化；
  \item 走访、就业联系和其他业务记录。
\end{itemize}

确认前的草稿不会进入正式工作留痕和首页统计。导出时可以保留来源集合、来源记录 ID、来源标题和来源状态；这是交接和审计的重要依据。

\uiShot{09-worklog-drafts.png}{工作留痕草稿：先核对来源，再确认归档}{fig:worklog-drafts}

\chapter{任务、通知、日程与工作节点}

\section{工作任务}

任务支持标题、说明、截止日期、优先级、状态、学生或班级关联、来源记录、附件和负责人。首页按逾期、今日、近期和待确认分组；任务详情可以直接跳回关联学生、谈话、课题或住宿批次。

完成任务不等于自动写入“已经发生”的工作事实。对需要留痕的任务，完成后先生成草稿，再由老师确认。

\uiShot{10-tasks.png}{工作任务：按日期、优先级和状态安排下一步工作}{fig:tasks}

\section{通知 AI 工作区}

通知页面采用“原文输入 / AI 结果”双栏结构。老师主动粘贴通知或导入文本，系统可以整理标题、来源、接收时间、适用对象、重点、截止时间、待办、需要核验的信息和原文证据片段。

\begin{figure}[H]
  \centering
  \includegraphics[width=\textwidth]{11-notice-ai.png}
  \caption{通知 AI：左侧保留原文，右侧生成可编辑、可核对的结果}
  \label{fig:notice-ai}
\end{figure}

结果支持编辑、重新生成、复制、保存为草稿、提取截止时间和适用对象。确认后可以分别转为任务或待确认工作留痕。系统不会后台读取钉钉、学校后台或其他浏览器通知，也不会把 AI 结果直接当作学校正式通知。

通知原文默认保存摘要和哈希；只有用户明确选择保存时才把完整原文放入本地资料库。外部来源需要人工触发，页面显示 URL、标题、抓取时间、最近核验时间和引用片段。

\uiShot{44-notice-ai.png}{通知处理：原文、重点和可转化动作放在同一工作区}{fig:notice-ai-new}

\section{工作节点分类}

工作节点分类是可配置字典，默认包含日常管理、社区管理、劳动卫生、班级建设、走访、就业、科研和材料报送等类型。自定义类别不会重命名旧记录；删除分类目录也不会删除已有工作事实。

\uiShot{17-activities-worklogs.png}{活动与工作留痕：记录活动过程、参与人和材料附件}{fig:activities-worklogs}

\chapter{课表、考勤、查课与随机点名}

\section{班级课表}

班级课表支持 CSV/XLSX 导入、周视图、列表视图、学期筛选、班级筛选和打印。课程包含班级、学期、星期、起止节次、周次、教师、教室和备注。课表是请假课时计算和查课页面的基础来源。

\uiShot{30-class-schedule.png}{班级课表：按学期、班级和周次查看课程安排}{fig:class-schedule}

\section{考勤与请销假}

请假记录包含学生、请假类型、起止时间、原因、审批状态、附件和销假状态。系统将请假时间与课表节次相交，计算覆盖课时；已批准请假、未记录考勤、旷课和迟到分开统计，缺失考勤不直接当作旷课。

\uiShot{11-leave-attendance.png}{请假与考勤：申请、审批、销假和课时统计分开记录}{fig:leave-attendance}

\section{查课看板}

查课页面显示今日课程、正在上课、未查课程、未到人数、迟到人数和异常统计。查课先记录发现的事实，再填写处理措施、跟进期限和附件。查课完成后可生成待确认工作留痕。

\uiShot{31-class-checks.png}{查课看板：今日课程、课堂状态和异常处理入口}{fig:class-checks}

\section{随机点名}

随机点名支持多班级选择、全选、清空、单人或多人抽取、不重复抽取、随机种子、结果复核和大屏准备状态。随机数由本地随机源决定，AI 不参与抽取。

点名结果在人工复核后保存为一次点名会话，包含日期、班级、参与学生、抽取数量、随机种子、模式和结果。重复保存同一结果不会重复生成留痕草稿；投屏前应确认教室和隐私环境。

\uiShot{32-roll-call.png}{随机点名：选择班级、设置人数、复核结果并保存会话}{fig:roll-call}

\chapter{住宿、走读与查寝}

\section{校外住宿与走读}

校外住宿是学生外宿、走读申请的独立入口，支持申请日期、起止日期、申请理由、紧急联系人、审批状态、附件和查询。它与宿舍楼栋排宿分开维护。

\begin{figure}[H]
  \centering
  \includegraphics[width=\textwidth]{09-dorm-special.png}
  \caption{住宿相关入口：校外住宿、走读记录与专项管理上下文}
  \label{fig:dorm}
\end{figure}

\section{楼栋、房间与住宿批次}

住宿专项包括楼栋、房间、床位、批次和学生排宿：

\begin{itemize}
  \item 楼栋：校区、名称、性别限制、启用状态和备注。
  \item 房间：所属楼栋、楼层、房间号、床位数、床位编号和房间状态。
  \item 批次：学年、学期、批次类型、状态和说明。
  \item 排宿：稳定学生 ID、学号快照、楼栋、房间、床位、入住日期和退宿日期。
\end{itemize}

排宿方案先生成预览，检查容量、性别限制、重复入住、床位冲突和未分配学生。只有人工确认后才更新学生当前楼栋、房间和床位快照。

\section{调宿与历史轨迹}

调宿记录保留原位置、新位置、原因、办理日期、操作人和附件。它不会覆盖原有住宿轨迹；学生档案显示当前快照，住宿历史显示每一次变化。

\section{查寝与异常闭环}

查寝记录保存批次、楼栋、房间、检查日期、检查人、结果和照片/附件。发现问题后单独建立异常记录，包含异常类别、等级、描述、处理期限、状态、处理结果和关闭时间。

住宿统计支持按楼栋、楼层、房间、月份和未处理异常筛选。查寝事实和异常处理状态分开，避免把“已经检查”误解为“问题已经解决”。

\uiShot{33-dorm-inspections.png}{宿舍查寝：按楼栋和房间记录检查结果，再单独处理异常}{fig:dorm-inspections}

\uiShot{34-dorm-committee.png}{住宿专项：楼栋、房间、床位和学生分配的集中入口}{fig:dorm-committee}

\chapter{班团组织、党团发展与量化考评}

\section{班团组织与班委}

班团组织用于维护班级组织、岗位、任期、代理、空缺、继任和学生任职经历。默认岗位包括班长、团支书、学习委员、生活委员和心理委员，也可以增加权益委员、就业委员、科研委员、宣传委员和文体委员。

岗位字典属于当前工作区配置，一人可以担任多个岗位。任职关系优先使用稳定学生 ID，保留学号快照和任期日期。

\uiShot{13-class-organizations.png}{班团组织：自定义岗位、任职关系和班级范围一处维护}{fig:class-organizations}

\section{班委考核与量化考评}

班委考核按周期录入评价日期、岗位、等级、备注和改进建议。等级由老师人工选择，AI 只能生成评价措辞草稿。

量化考评使用规则版本、评价维度、权重、上限和等级阈值。明细记录包含学生、学期、维度、加减分、来源、证据附件、备注和核验状态。支持 Excel 导入、手动加分/扣分、学期/年级/班级筛选、明细下钻、排名和导出。

基础分、累计加分、累计扣分、最终分和平均分由明细派生；没有记录显示“未记录”，不自动当作零。规则版本和导入来源保留，以便解释同一学生在不同周期的结果。

\uiShot{37-assessment.png}{量化考评：规则版本、积分明细、核验状态和排名入口}{fig:assessment}

\begin{figure}[H]
  \centering
  \includegraphics[width=\textwidth]{14-party-development.png}
  \caption{党团发展：流程节点、日期、材料和状态维护}
  \label{fig:party}
\end{figure}

\section{党员发展与团员发展}

党员发展按规则版本维护申请、培养、考察、接收预备党员和转正等阶段、日期、材料、来源和状态。团员发展采用国家基线加学校附加节点，核心节点不能删除。

系统只做流程提醒、日期计算和材料完整性检查，不代替党团组织的政治审查、资格判断、审批或组织结论。日期缺失保持为空，不自动填成当天。

\uiShot{15-league-development.png}{团员发展：学校基线节点与自定义节点并列维护}{fig:league-development}

\chapter{成绩、学业预警、帮扶与班级分析}

\section{成绩管理}

成绩可以按学期、课程、学生和班级导入与查看。页面提供单学生成绩趋势、课程明细、挂科课程、挂科门数、GPA、班级均值和学期对比。缺失成绩显示“未记录”，不直接当作零分。

\begin{figure}[H]
  \centering
  \includegraphics[width=\textwidth]{07-grades-support.png}
  \caption{成绩管理：趋势图、课程风险和学业帮扶联动}
  \label{fig:grades}
\end{figure}

\section{学业预警与帮扶}

预警记录包含来源课程、风险等级、触发原因、解除状态、核查日期和备注。帮扶记录包含目标、措施、负责人、周期、回访和附件。识别出风险后，可以从成绩页面直接创建帮扶任务或回到学生档案补充谈话记录。

预警辅助分析可以生成解释草稿，但不能自动修改预警结论、帮扶状态或学生事实。学业帮扶的完成要有来源、日期和人工确认。

\uiShot{18-academic-warning.png}{学业预警：风险等级、触发课程和处理状态清楚分列}{fig:academic-warning}

\uiShot{19-academic-support.png}{学业帮扶：目标、措施、负责人和回访计划形成连续记录}{fig:academic-support}

\section{班级综合分析}

班级分析支持班级、年级、学期和时间范围筛选，聚合学生人数、成绩、GPA、挂科、考勤、谈话、重点关注、预警、资助、获奖、处分、活动、班委履职、就业和科研等指标。

敏感指标默认只展示聚合数量，进入个人明细前需要访问控制。缺失记录显示“未记录”。分析结果是派生计算或缓存，不新增重复事实表；比较两个学期时保留统计口径和过滤条件。

\uiShot{20-grades-analysis.png}{成绩分析：趋势、挂科课程和班级对比支持下钻}{fig:grades-analysis}

\uiShot{42-academic-class-analysis.png}{班级学业分析：从班级汇总进入课程和学生明细}{fig:academic-class-analysis}

\uiShot{43-class-summary.png}{班级综合摘要：关键指标、学期对比和下一步工作入口}{fig:class-summary}

\chapter{假期、考勤、奖惩、资助与就业}

\section{假期去向与请销假}

假期去向记录学生假期所在地点、联系人、出行方式、起止日期、风险提示和确认状态。请销假记录关注审批、出入时间、销假和附件。两个模块分别保存“计划去向”和“实际请假流程”，避免相互覆盖。

\uiShot{12-approvals-attendance.png}{请销假管理：审批状态、销假动作和历史记录一目了然}{fig:leave-approval}

\section{评优、奖惩与证书}

评优榜样和奖惩管理支持荣誉、处分、等级、日期、来源、状态、说明和附件。证书图片、PDF 和常见文档进入附件仓，业务记录只保存附件 ID 和文件哈希。

证书识别只生成学生、证书名称、等级和日期草稿。老师需要选择稳定学生、检查字段并人工确认后，结果才写入正式奖惩记录。当前不承诺任意 PDF 的自动视觉识别。

\uiShot{16-rewards-grants.png}{奖惩管理：荣誉、处分、证书附件和批量维护入口}{fig:rewards}

\section{奖助勤与伙食补助}

奖助勤模块按学期维护奖学金、助学金、勤工助学、困难等级、伙食补助、认定状态、发放状态、材料和来源。困难认定材料可以使用附件 ID 归档，导出前显示敏感字段和范围确认。

资助数据不能被 AI 自动认定、修改或生成结论。班级分析默认只显示聚合结果。

\uiShot{21-grants-and-assistance.png}{奖助勤补：项目、金额、批次和发放状态集中管理}{fig:grants}

\section{毕业生就业}

就业模块支持就业状态、就业意向、单位、岗位、联系记录、求职导航、简章摘要、材料清单和通知草稿。学生就业信息优先关联稳定学生 ID，旧的自由文本保留为兼容快照，不静默猜测学生身份。

就业防骗台账保存单位名称、风险等级、理由、来源 URL、核验时间和备注。外部链接必须经过 HTTPS、私网地址、凭据和危险协议检查，打开外链由用户明确触发。

\uiShot{24-employment.png}{就业工作区：就业状态、意向、联系记录和资源入口统一查看}{fig:employment}

\uiShot{39-employment-safety.png}{就业防骗：风险等级、理由、来源和最近核验时间清楚记录}{fig:employment-safety}

\section{业务档案、重点学生与心理摸排}

业务档案用于承载学校自定义的专项记录；重点学生档案用于维护重点关注类型、关注等级、危机状态、发现方式、跟进频率和解除状态。两者都应绑定稳定学生 ID，并通过学生档案、时间线和首页待办回看，不把专题记录复制成另一份学生主体。

心理摸排支持按班级、年级和关注状态形成工作清单，记录兴趣、人工重点关注、谈话主题、后续建议、访问范围和修改/查看审计。它用于组织辅导员的日常跟进，不是心理测评诊断工具；群体展示优先使用脱敏聚合，个人详情必须经过受控访问。

\uiShot{22-focus-students.png}{重点学生：按关注类型、等级和跟进状态整理工作清单}{fig:focus-students}

\uiShot{23-psych-screening.png}{心理摸排：用受控范围记录兴趣、主题和后续工作安排}{fig:psych-screening}

\chapter{活动、科研、竞赛、资料与模板}

\section{活动开展与批量图片}

活动记录包含活动、日期、学期、班级、组织者、参与学生、摘要、状态和附件。活动照片支持批量选择、哈希去重、按活动/日期/班级/学生筛选和导出。

批量上传遵循“附件先写入、业务记录后提交、全部成功才确认”的边界。业务记录保存失败时清理没有其他引用的附件；清理失败会进入诊断字段，不把孤立附件伪装成完成。

\uiShot{17-activities-worklogs.png}{活动留痕：参与人、活动摘要和图片材料集中管理}{fig:activity-records}

\section{科研课题}

科研课题按申请书准备、申报、评审/立项、开题、研究中、中期检查、结题材料和已结题管理。每个阶段记录截止日期、任务、附件、来源和变更时间线。同一项目、同一阶段重复点击只生成一条有效阶段任务；进入下一阶段后才创建新的阶段任务。

\uiShot{36-research-projects.png}{科研课题：按阶段查看项目、截止日期和阶段任务}{fig:research-projects}

\section{科创竞赛}

竞赛资源支持分类、竞赛名称、主办方、官网、报名地址、截止日期、来源和核验状态。竞赛报名记录关联学生、项目、角色、分工、报名状态、获奖等级、附件和备注。

旧的自由文本参赛信息继续保留，不把文本姓名自动猜测成学生身份。模板材料与附件可以被项目引用，但删除一条报名记录时不会误删共享附件。

\uiShot{38-employment-competition.png}{竞赛资源：报名项目、成员分工、截止时间和材料入口}{fig:competition}

\section{政策智库、素材收藏与模板库}

政策智库用于政策文件、制度和流程资料；素材收藏用于讲话稿、海报、材料和外部参考；模板库用于 CSV/XLSX 模板、Word 模板、工作包和表单映射。每条资料可以记录分类、来源、版本、更新时间、核验状态、关键词和适用场景。

\section{学习助手}

学习助手把政策阅读、日常材料学习和个人工作笔记放在同一入口。资料可以记录类型、来源、标签、摘要、重点、学习状态和进度；笔记单独保存并关联资料，不会覆盖原始政策或模板。适合把“先看材料、再形成自己的工作要点”变成可回看的过程。

外部资料只在用户明确打开或导入时进入工作区。涉及学校制度的内容应保留来源、发布日期和版本，并以学校正式文件为准；AI 可以帮助提炼摘要和检查缺失字段，但不能把非正式参考材料自动当作制度结论。

\uiShot{28-learning-assistant.png}{学习助手：资料进度、笔记和摘要形成个人工作资料库}{fig:learning-assistant}

\uiShot{25-policy-library.png}{政策资料：按来源、类别和适用场景查找工作依据}{fig:policy-library}

\uiShot{26-material-library.png}{素材收藏：讲话稿、案例和日常材料随时检索}{fig:material-library}

\section{一生一表与 Word 反向汇总}

一生一表支持内置版本化 Word 模板和学校自定义模板。受支持模板使用固定占位符。生成前显示学生范围、字段映射和缺失字段；生成后保留任务和模板版本。
\begin{itemize}
  \item \placeholder{\{\{student.full\_name\}\}}
  \item \placeholder{\{\{student.student\_number\}\}}
  \item \placeholder{\{\{student.class\_name\}\}}
\end{itemize}

Word 反向汇总只支持使用受支持占位符或内容控件的模板。导入前显示字段映射和与现有数据的差异，冲突字段不静默覆盖，结果默认导出结构化 CSV，不自动写回学生事实。

\uiShot{27-templates-learning.png}{模板中心：模板版本、适用场景和生成任务统一管理}{fig:template-center}

\uiShot{49-content-and-forms.png}{内容与表单：学校可以发布资料并按范围查看已读状态}{fig:content-forms}

\chapter{实用工具与资料入口}

\section{工作通讯录}

工作通讯录用于保存学院、学校部门、校外协作单位和日常办公联系人。每条联系人可以维护姓名、部门、岗位、联系方式、适用场景、备注和最近核验时间；联系方式属于敏感信息，查看与导出遵循访问锁和最小范围原则。

工作通讯录与家长联系严格分离，也不会根据姓名或电话自动把联系人关联为学生、家长或业务事实。需要与学生关联时，应在具体业务记录中明确选择学生并保存当次联系的关系快照。

\uiShot{29-contacts-schedules.png}{工作通讯录：部门联系人、办公方式和适用场景快速检索}{fig:contacts}

\section{工具入口}

工具入口集中维护学校智慧学工、教务、就业、住宿、日常材料、竞赛和自定义网站。链接记录名称、分类、URL、说明、排序、收藏、核验状态和更新时间。系统只允许经过校验的 HTTPS 外链，拒绝危险协议、带凭据 URL 和明显私网地址。

\begin{figure}[H]
  \centering
  \includegraphics[width=\textwidth]{10-tools.png}
  \caption{工具入口：按分类维护学校办公网址和常用工具}
  \label{fig:tools}
\end{figure}

\section{抽签、分组和轮值}

实用工具默认只在当前会话中处理名单和结果，不保存名单。名单来源可以是手动输入、当前班级或当前筛选学生。只有明确点击“保存为任务”“保存为留痕”或“保存为名单”后才进入业务数据。

\begin{itemize}
  \item 抽签：支持单人/多人、不重复抽取、本地随机源、结果摘要和随机种子复核。
  \item 分组：按组数或每组人数生成分组，可选择基础均衡条件。
  \item 轮值：按起始日期、周期、人员和轮次数量生成任务草稿。
  \item AI 辅助：可以生成分组说明或通知文字，但不参与随机数决定，不自动修改学生关系。
\end{itemize}

\section{日期计算与名单清理}

日期计算支持回访日期、截止日期、阶段间隔、教学周期和倒计时。名单清理支持拆分、去重、排序、空白清理、格式统一和 CSV 导出。默认不把临时清理结果写入学生台账。

\uiShot{40-tools-utilities.png}{实用工具：抽签、分组、轮值、日期计算和名单清理集中使用}{fig:utilities}

\uiShot{41-utilities-draw.png}{随机抽签：名单、抽取人数和不重复结果在一次操作中完成}{fig:utility-draw}

\chapter{AI 助手：在需要的地方帮你完成文字工作}

\section{AI 的正确定位}

AI 适合处理整理、提炼和起草工作：把一段通知整理成重点，把一名学生的近期记录整理成 briefing，把分析结果写成一份易读的说明。老师仍然掌握最终判断和保存按钮，AI 生成的内容都可以先修改、再确认。

\section{AI 入口覆盖}

AI 可以在以下具体场景中使用：

\small
\begin{longtable}{>{\bfseries}p{31mm} p{65mm} p{47mm}}
\caption{AI 页面级入口与输出边界}\\
\toprule
入口 & 可生成内容 & 不会自动执行的动作 \\
\midrule
学生档案 & 学生摘要、成绩趋势解释、资料缺失提醒 & 不改学生事实和历史记录 \\
谈心谈话 & 会前 briefing、谈话草稿、回访建议 & 不生成心理诊断或危机结论 \\
通知工作区 & 重点提取、截止日期、适用对象、通知草稿 & 不后台读取外部通知，不直接发通知 \\
成绩与预警 & 风险解释、帮扶建议草稿 & 不改变成绩、预警或帮扶状态 \\
就业与竞赛 & 简章摘要、材料清单、通知草稿 & 不认定单位安全，不替学生报名 \\
科研与模板 & 阶段清单、字段缺失检查、表单映射 & 不审批课题，不静默覆盖冲突 \\
证书识别 & 证书字段草稿 & 不自动认定奖项，不自动绑定学生 \\
住宿查课查寝 & 冲突解释、异常摘要、材料提示 & 不决定床位、不关闭异常、不确认考勤 \\
同步冲突 & 差异解释 & 不按最后修改时间自动覆盖同字段 \\
\bottomrule
\end{longtable}
\normalsize

\uiShot{45-ai-workbench.png}{AI 工作台：模型状态、当前页面建议、来源和待确认草稿集中查看}{fig:ai-workbench}

\section{模型连接与调用前检查}

在 AI 设置中添加模型连接，填写连接名称、协议、模型、用途和额度。连接状态必须区分“未配置”“已配置但不可调用”和“可调用”。文本模型可用不代表语音转写可用，语音需要单独的转写模型和服务商音频端点。

请求发送前，系统至少检查模型连接、用途是否登记、当次授权、敏感字段范围、请求大小、图片模型能力和必要的转写模型。模型 API Key 只保存在桌面端安全存储或当前浏览器会话，不进入业务备份、交换包和普通日志。

\section{敏感字段授权}

学生身份、联系方式、心理、纪律、资助、预警、重点关注、家庭信息和附件属于敏感分类。生成前可以预览本次即将发送的字段和来源；敏感分类必须逐次明确授权，取消授权不会发出请求。

发送给模型的内容会去掉工作台内部编号，页面会把本次使用的范围和来源列出来，方便你在提交前快速核对。

\section{建议中心与人工确认}

AI 结果先进入建议或草稿中心，状态包括草稿、待审核、已查看、已采纳和已驳回。老师可以编辑、反馈、重试或转为任务、谈话和工作留痕。来源被修改或删除后，建议会显示“来源已变化”或“来源已删除”，不能直接作为最新事实确认。

模型请求失败、取消、超时、响应过大或返回格式错误时，输入和重试范围保留，原始响应不写入普通业务记录。成功提示必须等待真实持久化完成；失败时恢复页面状态并保留重试按钮。

\section{心理工作与语音转写}

语音转写需要在模型设置中配置支持音频的服务。每次发送前重新确认授权，原始录音默认不保存；模型服务未配置时，可以先查看页面和整理流程，不会误把录音当作已完成记录。

\section{使用 AI 时记住一件事}

涉及学生事实、住宿安排、心理与危机、纪律处分、资助认定、预警、奖惩、科研审批、党团发展和考勤的内容，最后都由老师核对并点击确认。AI 负责整理和起草，不替你做决定。

\chapter{资料、附件、导入导出与模板}

\section{附件仓}

照片、证书、活动图片、查课查寝材料、困难认定材料和科研文件进入附件仓。业务记录只保存附件 ID、文件名、大小、哈希和引用关系。附件支持缩略图、预览、下载、删除、索引导出和重复识别。

删除附件前检查共享引用；仍被其他记录引用的附件不会直接清理。业务记录写入失败时，刚写入但没有其他引用的附件回滚；清理失败进入诊断状态，避免留下不可追溯的“已完成”提示。

\section{导出范围与敏感提示}

导出前显示范围、字段、用途和敏感提醒。学生、照片、家庭、心理、资助、纪律、就业、交换包和备份等导出需要确认保存位置和使用范围。导出文件不包含模型 API Key、内部授权令牌、设备令牌或学生数据之外的运行时凭据。

\section{打印与工作包}

宿舍名单、学生名册、课表、班级分析、表单任务和模板材料可以按当前筛选结果打印或导出。打印前应核对页眉、筛选范围、日期和敏感字段；不要把包含完整联系方式、心理或资助明细的工作包放在公共打印区。

\chapter{备份、恢复、迁移与跨端互通}

\section{存储层次}

\begin{table}[H]
\centering
\caption{不同运行方式的存储与互通能力}
\begin{tabularx}{\textwidth}{B{35mm} P{40mm} Y}
\toprule
运行方式 & 主要存储 & 使用边界 \\
\midrule
浏览器 & IndexedDB、本地加密快照 & 关闭浏览器期间不能后台写文件；支持手动导出和用户选择目录保存。 \\
离线 HTML & 浏览器本地工作区 & 适合无网络演示和临时使用；新版本需要重新获取离线文件。 \\
Electron & SQLite、WAL、加密附件仓、系统安全存储 & 适合长期单机使用，支持数据目录、恢复口令和定时备份。 \\
局域网主机 & Electron 主机的 SQLite 与附件仓 & 手机、网页和其他桌面端通过 HTTPS、配对令牌和同步队列接入。 \\
\bottomrule
\end{tabularx}
\end{table}

\section{备份策略}

桌面端在工作台运行期间按定时器检查备份条件，通常按每日一次或变更量阈值生成加密备份。浏览器端保留本地加密快照；只有在用户明确选择文件夹并且浏览器支持 File System Access API 时，才会写入用户指定目录。浏览器关闭后无法后台写文件，页面必须明确显示这一限制。

\begin{figure}[H]
  \centering
  \includegraphics[width=\textwidth]{08-backup-migration.png}
  \caption{备份与迁移：加密保险库、恢复、交换包和数据目录入口}
  \label{fig:backup}
\end{figure}

\section{恢复顺序}

恢复前先生成当前工作区回退快照，再做完整性检查、差异预览和附件索引校验。恢复失败自动回滚，不覆盖当前有效数据。桌面端的恢复口令用于换机恢复工作区主密钥，口令只显示一次；忘记恢复口令不能通过后门解锁。

遇到密钥损坏、SQLite/WAL 损坏、附件解密失败或仓储校验失败时，应停止继续写入，保留原数据目录和恢复包，进入“数据修复与恢复”页面。不要先卸载软件，更不要手动删除用户数据目录。

\uiShot{46-backup.png}{备份与恢复：备份状态、恢复点和数据修复入口集中显示}{fig:backup-new}

\section{局域网数据中枢}

局域网主机由 Electron 主进程提供 SQLite、加密附件仓和同步服务。客户端不直接打开同一个 SQLite 文件，而是通过 HTTPS、证书指纹和一次性配对码或二维码接入主机。

同步操作包含工作区、设备、操作 ID、幂等键、集合、记录、基础修订号和字段补丁。不同字段修改可以自动合并；同字段冲突保留双方值、设备、时间和来源，进入冲突收件箱，人工选择主机值、客户端值或手动编辑。

附件通过独立的分块上传接口传输，支持哈希校验、断点续传和失败清理；附件二进制不放入业务同步记录。设备可以暂停、恢复和撤销，配对、同步、冲突和附件操作写入本地最小审计。

\section{三端互通的实际流程}

\begin{enumerate}
  \item 在 Electron 主机启动局域网数据中枢，确认 HTTPS 地址和证书指纹。
  \item 在主机生成一次性配对码或二维码，核对有效期和工作区标识。
  \item 手机或网页端输入/扫描配对信息，人工核对地址和指纹后向主机发起请求。
  \item 主机确认设备，客户端取得一次性设备令牌并连接。
  \item 客户端先提交离线队列，再拉取远程修订；同步完成后刷新冲突收件箱。
  \item 遇到同字段冲突时人工选择，解决结果形成新的修订，其他设备再次拉取。
  \item 主机定期备份工作区；换机时使用恢复包恢复主机，再重新配对客户端。
\end{enumerate}

主机离线时，客户端可以保留加密离线队列；浏览器关闭期间不会继续后台同步。局域网首期只面向同一网络，不等同于公网协作平台，也不自动信任二维码中的证书。

\uiShot{47-lan-sync.png}{局域网同步：主机、设备、离线队列和冲突收件箱一屏查看}{fig:lan-sync}

\uiShot{48-platform-bridge.png}{跨端数据中心：备份、交换包、数据目录和恢复操作统一管理}{fig:platform-bridge}

\chapter{隐私、安全与访问审计}

\section{本地优先不等于自动安全}

本地优先意味着数据默认留在当前设备或明确配置的主机上，并不表示设备没有病毒、没有备份泄露风险或浏览器一定具备数据库级加密。应使用操作系统账户、磁盘加密、屏幕锁、受控目录和独立恢复口令保护设备。

\section{访问锁与最小范围}

家长电话、完整联系方式、心理、危机、资助、纪律、就业和重点关注等信息默认最小化展示。查看敏感详情或导出时，页面显示范围和字段，并记录访问审计。班级领导统计默认使用聚合结果，个人明细需要重新进入受控档案。

\section{审计记录}

访问审计记录访问、导出、敏感授权、AI 请求、建议确认、同步配对、设备撤销、冲突解决、恢复和附件操作的最小元数据。普通日志不保存 API Key、恢复口令、设备令牌、完整 AI 请求、完整模型响应、原始录音或学生附件二进制。

\section{来源链与可追溯}

工作留痕草稿、AI 建议和通知结果保留来源集合、来源记录、来源哈希、来源状态和最近核对时间。来源发生变化时必须重新核对，来源删除后只能重新创建。这样可以区分“当时生成的建议”和“当前仍然有效的事实”。

\section{学校制度优先}

以下事项必须以学校制度、专业人员和正式系统为准：心理危机处置、医学判断、纪律处分、资助认定、学籍变更、党团审批、住宿正式分配、成绩认定和就业统计上报。工作台的记录、提醒和分析不能替代这些流程。

\uiShot{51-audit-and-recovery.png}{访问审计与恢复：查看敏感访问、导出和恢复操作记录}{fig:audit-recovery}

\chapter{平台差异、故障处理与长期维护}

\section{浏览器、离线 HTML 和 Electron 的选择}

浏览器适合快速开始和单机办公；离线 HTML 适合无网络环境下的演示或临时工作；Electron 更适合长期使用、附件管理、数据目录、定时备份和局域网主机。涉及学生照片、证书和大量记录时，优先使用 Electron 或受控主机，并建立可恢复备份。

\section{常见问题处理}

\small
\begin{longtable}{>{\bfseries}p{45mm} p{94mm}}
\caption{常见问题与建议动作}\\
\toprule
现象 & 建议动作 \\
\midrule
导入三列成绩被识别为新增学生 & 检查学号是否完整，确认导入模式为合并更新；查看预览中的匹配列和待核对清单，不按姓名手工猜测。 \\
保存后刷新看不到数据 & 停止继续操作，先查看存储状态和错误提示，保留原文件并点击重试；不要用导出文件覆盖原工作区。 \\
照片或证书显示不完整 & 检查附件仓、附件 ID 和引用状态，先保留原目录与备份，再重新读取索引。 \\
AI 按钮不能使用 & 先检查模型连接、用途授权、敏感字段授权、模型类型和额度；文本可用不代表语音端点可用。 \\
手机无法连接主机 & 使用电脑局域网地址，不要使用手机上的 \code{127.0.0.1}；核对 HTTPS、证书指纹、配对码有效期和主机确认状态。 \\
同步出现冲突 & 进入冲突收件箱，比较双方值和来源后人工选择；不要关闭页面，也不要以最后修改时间替代判断。 \\
升级后看不到数据 & 停止写入，保留旧数据目录和恢复点，先做健康检查，再按恢复向导处理；不要先卸载或删除目录。 \\
浏览器关闭后没有自动备份 & 这是浏览器能力边界；使用桌面端定时备份，或在关闭前手动导出加密备份。 \\
\bottomrule
\end{longtable}
\normalsize

\section{长期维护的推荐节奏}

\begin{enumerate}
  \item 开学建档：先导入基础身份和班级，再补照片、家长、导师、班主任和住宿字段。
  \item 日常工作：每条谈话、联系、查课、查寝和活动记录都写清下一步和日期。
  \item 期中期末：使用学号或稳定 ID 增量导入成绩和考勤，先看预览，再确认写入。
  \item 月度复盘：查看首页逾期项、未回访、预警、帮扶、住宿异常和待确认草稿。
  \item 关键操作前：批量删除、覆盖导入、迁移数据目录、更新程序和恢复工作区前先生成恢复点。
  \item 交接：导出必要的工作包和来源说明，清理不再需要的临时文件，不把完整敏感明细发到公共群组。
\end{enumerate}

\chapter{功能索引与数据原则}

\section{功能索引}

\small
\begin{longtable}{>{\bfseries}p{32mm} p{46mm} p{59mm}}
\caption{全量模块索引}\\
\toprule
模块 & 主要工作 & 关键数据原则 \\
\midrule
今日概览 & 今日任务、待回访、草稿和系统状态 & 只显示摘要和最紧急事项 \\
工作任务 & 新建、关联、完成、回访和转留痕 & 来源明确，完成后人工确认 \\
学生台账 & 导入、筛选、批量编辑、照片和档案 & 稳定学生 ID 优先，姓名不自动匹配 \\
谈心谈话 & 记录、筛选、briefing、回访 & 不替代心理诊断，草稿需确认 \\
工作留痕 & 事实记录、查课查寝、来源回链 & 来源变更后重新核对 \\
校外住宿 & 走读申请、审批、联系人 & 与宿舍排宿分开 \\
假期去向 & 去向、联系人、起止日期和确认 & 计划与实际请假分开 \\
评优榜样/奖惩 & 荣誉、处分、证书和附件 & AI 不能自动认定结论 \\
请销假/考勤 & 请假、销假、迟到、旷课和课时 & 缺失不等于旷课 \\
工作节点 & 社区、劳动卫生、就业、科研等节点 & 分类可配置，旧记录不重命名 \\
班团组织 & 岗位、任期、代理、空缺和继任 & 一人多职，使用稳定 ID \\
党员/团员发展 & 阶段、日期、材料和提醒 & 不替代组织审批 \\
成绩/预警/帮扶 & 成绩趋势、风险、计划和回访 & 缺失显示未记录，不当零 \\
奖助勤补 & 认定、发放、伙食补助和材料 & 敏感字段最小化展示 \\
业务档案/重点学生 & 专题记录和受控详情 & 访问审计与聚合展示 \\
毕业生就业 & 意向、联系、资源和防骗 & 外链核验，不猜身份 \\
政策/素材/模板 & 文件、政策、讲话稿和模板 & 来源、版本和核验状态可追溯 \\
科创竞赛/科研 & 资源、报名、项目、阶段和附件 & 阶段任务去重，附件 ID 引用 \\
通讯录/课表 & 办公联系人、课程和打印 & 通讯录不自动关联学生 \\
工具/实用工具 & 外链、抽签、分组、轮值、日期、名单 & 临时结果默认不保存 \\
AI 智能工作台 & 建议、草稿、来源和审计 & 不自动修改高风险事实 \\
平台联动/备份 & 交换包、恢复、局域网和迁移 & 先验证再提交，失败回滚 \\
访问审计 & 查看、导出、AI、同步和恢复记录 & 不记录密钥和完整原文 \\
\bottomrule
\end{longtable}
\normalsize

\section{最重要的五条原则}

\begin{enumerate}
  \item 用稳定 \code{student\_id} 关联学生业务，学号只做当前和历史快照。
  \item 导入默认合并更新，空白保留原值；清空、覆盖、删除必须明确确认。
  \item 所有来源转换先生成草稿，人工确认后才进入正式事实或工作留痕。
  \item 附件只通过附件 ID 进入业务记录，删除时检查共享引用并保留审计。
  \item AI、同步和恢复都不能绕过权限、来源核对、人工确认和真实持久化。
\end{enumerate}

\chapter{开始使用前的三项准备}

\section{先熟悉，再录入}

第一次打开时，先用示例数据浏览首页、学生台账、成绩、谈话、住宿和资料模块。熟悉筛选、编辑、保存、导出和返回路径后，再清理示例记录并录入真实数据。

\section{从最短路径开始}

建议按下面的顺序建立自己的工作区：

\begin{enumerate}
  \item 在“设置”中填写称呼、学院和常用偏好。
  \item 在学生台账导入基础名单，先核对学号、姓名和班级数量。
  \item 生成第一份加密备份，再补充照片、家长、导师、班主任和住宿信息。
  \item 日常从“今日概览”处理任务、回访和待确认记录，月底按班级和学期做一次分析。
\end{enumerate}

\section{把反馈变成下一步}

遇到问题时，记录“所在模块、筛选条件、点击的按钮、页面提示和刷新后的结果”，再把不含学生资料的截图发给维护人员。建议功能最好同时说明使用场景和期望结果，这样更容易快速形成真正省时间的改进。

\begin{warningbox}
本手册的截图和示例数据仅用于演示。录入真实资料前，请先用少量脱敏数据完成导入、保存、回读和备份；不要把学生资料、备份口令或模型密钥发送到公共群组。
\end{warningbox}

\enlargethispage{4\baselineskip}
\section{反馈与持续改进}

反馈问题时请描述“进入哪个模块、选择了什么范围、点击了什么按钮、页面提示什么、刷新后数据如何变化”，并提供不含学生信息的截图或虚构样例。功能建议应说明实际工作场景、当前耗时、预期结果和涉及的数据范围，这比只描述“增加一个按钮”更容易形成可用的业务闭环。

\end{document}
